\section{Principios de Seguridad}

El sistema implementa seguridad en múltiples capas, siguiendo el principio de \textbf{defense in depth}. Todas las validaciones críticas se ejecutan en el backend (base de datos), no en el frontend.

\section{Autenticación}

La autenticación se gestiona mediante \textbf{Supabase Auth}, que proporciona:

\begin{itemize}
    \item Registro e inicio de sesión con email y contraseña.
    \item Contraseñas cifradas con bcrypt (gestionado por Supabase).
    \item Manejo seguro de sesiones mediante JSON Web Tokens (JWT).
    \item Renovación automática de tokens.
    \item Protección contra ataques de fuerza bruta (rate limiting).
\end{itemize}

\subsection{Login Dual (CI / Email)}

El sistema permite el inicio de sesión mediante CI (Carnet de Identidad) o correo electrónico:

\begin{enumerate}
    \item El estudiante ingresa su identificador (CI o email) y contraseña.
    \item Si el identificador contiene ``@'', se trata como email y se autentica directamente.
    \item Si no contiene ``@'', se busca el CI en la tabla \texttt{student} para obtener el email asociado.
    \item Se autentica contra Supabase Auth con ese email y la contraseña proporcionada.
\end{enumerate}

\section{Control de Acceso (RLS)}

\textbf{Row Level Security} (RLS) es el mecanismo principal de autorización. Cada tabla tiene políticas que definen qué filas puede leer, insertar, actualizar o eliminar cada usuario según su rol.

\subsection{Políticas para Estudiantes}

\begin{itemize}
    \item \textbf{student}: SELECT y UPDATE solo de su propio registro. INSERT permitido (registro público).
    \item \textbf{exam}: SELECT de sus propios exámenes. INSERT con \texttt{student\_id = auth.uid()}.
    \item \textbf{exam\_question}: SELECT de preguntas de sus propios exámenes.
    \item \textbf{student\_answer}: INSERT y UPDATE solo en exámenes en progreso.
    \item \textbf{level}: SELECT de todos los niveles (para visualizar resultado).
    \item \textbf{question}: SELECT de todas las preguntas (para responder el examen).
\end{itemize}

\subsection{Políticas para Administradores}

\begin{itemize}
    \item \textbf{admin}: SELECT de su propio registro.
    \item \textbf{level, question, question\_option}: CRUD completo.
    \item \textbf{exam}: SELECT de todos los exámenes (para reportes y dashboard).
    \item \textbf{exam\_config}: SELECT y UPDATE.
    \item \textbf{audit\_log}: SELECT de todos los registros de auditoría.
\end{itemize}

\section{Protección de Datos Históricos}

Una de las reglas de negocio más críticas establece que \textbf{los resultados históricos nunca pueden modificarse}. Esto se implementa mediante un trigger SQL:

\begin{lstlisting}[language=PostgreSQL, caption=Trigger de protección de históricos]
CREATE OR REPLACE FUNCTION fn_prevent_historical_change()
RETURNS TRIGGER AS $$
BEGIN
    IF OLD.exam_id IS NOT NULL AND EXISTS (
        SELECT 1 FROM exam WHERE id = OLD.exam_id AND status = 'completed'
    ) THEN
        RAISE EXCEPTION 'No se pueden modificar respuestas de un examen completado'
            USING HINT = 'Los resultados históricos nunca pueden modificarse';
    END IF;
    RETURN NEW;
END;
$$ LANGUAGE plpgsql;

CREATE TRIGGER trg_prevent_historical_change
    BEFORE UPDATE ON student_answer
    FOR EACH ROW
    EXECUTE FUNCTION fn_prevent_historical_change();
\end{lstlisting}

\section{Auditoría}

Todas las acciones administrativas quedan registradas en la tabla \texttt{audit\_log}. Cada registro incluye:

\begin{itemize}
    \item Administrador que realizó la acción.
    \item Tipo de acción (CREATE, UPDATE, DELETE).
    \item Tabla y registro afectado.
    \item Detalles en formato JSONB.
    \item Fecha y hora exacta.
\end{itemize}

\section{Validaciones en Base de Datos}

Además de RLS y triggers, se utilizan constraints de base de datos para garantizar la integridad:

\begin{itemize}
    \item \textbf{CHECK constraints}: Validación de rangos de puntaje, límites de tiempo, valores positivos.
    \item \textbf{NOT NULL}: Campos obligatorios.
    \item \textbf{UNIQUE}: CI y email únicos en \texttt{student}, combinaciones únicas en tablas relacionales.
    \item \textbf{FOREIGN KEY}: Integridad referencial entre todas las tablas relacionadas.
    \item \textbf{ON DELETE RESTRICT}: Evita eliminaciones que dejarían datos huérfanos.
\end{itemize}

\section{Protección contra Ataques Comunes}

\begin{table}[H]
\centering
\begin{tabularx}{\textwidth}{l Y}
\toprule
\textbf{Amenaza} & \textbf{Mitigación} \\
\midrule
SQL Injection & Consultas parametrizadas mediante Supabase SDK \\
XSS & React escapea automáticamente el contenido HTML \\
CSRF & Tokens JWT en cada solicitud \\
Fuerza Bruta & Rate limiting de Supabase Auth \\
Acceso no autorizado a datos & RLS policies a nivel de fila \\
Modificación de datos históricos & Trigger SQL de protección \\
\bottomrule
\end{tabularx}
\caption{Mitigación de amenazas de seguridad}
\end{table}
